\documentclass[10pt]{article}

\usepackage[utf8]{inputenc}

\usepackage[T1]{fontenc}

\usepackage{amsmath}

\usepackage{amsfonts}

\usepackage{amssymb}

\usepackage[version=4]{mhchem}

\usepackage{stmaryrd}



\title{III }



\author{}

\date{}





\DeclareUnicodeCharacter{0131}{$\imath$}



\begin{document}

\maketitle

IIIАЛЯ ТЕОРЕМА - 1) III. т.: если $A, B, C$-три произвольные точки оси, то $\overline{A B}+\overline{B C}=\overline{A C}$, где $\overline{A B}$, $\overline{B C}, \overline{A C}$-длины направленных отрезков. ІІІ. т. обобщается для случая площадей ориентированных треугольников и объемов ориентированных тетраэдров (см. [1]).



\begin{enumerate}

  \setcounter{enumi}{1}

  \item III. т.: движение 1-го рода (не меняющее ориентации), отличное от поворота и переноса, является произведением переноса на поворот, ось к-рого параллельна направлению переноса (т. н. винтовое движение). Теорема доказана М. ШІалем (M. Chasles, 1830).

\end{enumerate}



М., $\underset{1969 .}{ }$ [1] М о д е н о в П1. С., Аналитическая геометрия,



III АIКА, $k$ - II $\mathbf{a}$ П $\mathbf{\kappa} \mathbf{a},-$ множество $k$ точек конечного проективного пространства $P(n, q)$, никакие три из к-рых неколлинеарны. Две Ші. считают әквивалентными, если существует коллинеация пространства $P(n, q)$, переводящая одну из них в другую. Нахождение максимального числа $m(n, q)$ точек ШII. в $\boldsymbol{P}(n, q)$, построение и классификация $m(n, q)$-ШІ. составляет главный вопрос при изучении ШІ., не решенный полностью (1984). Известны следующие результаты (см. [2], [3]):



$m(n, 2)=2^{n}: 2^{n}$-III. единственна (с точностью до эквивалентности) и является множеством точек $P(n, 2)$, не лежащих в фиксированной гиперплоскости; $m(2, q)=q+1, q$ - нечетное: $(q+1)$-III. в $P G(2, q)$ является коникой и едивственна;



$m(2, q)=q+2, \quad q$ - четное: $(q+2)$-III. в $\boldsymbol{P}(2, q)$, вообще говоря, не едивственна; $m(3, q)=q^{2}+1, q-$ нечетное: $\left(q^{2}+1\right)$-ШI. в $P(3, q)$ является эллиптич. квадрикой и единственна;



$m(3, q)=q^{2}+1, q$ - четное: $\left(q^{2}+1\right)$-III. в $P(3, q)$, вообще говоря, не единственна;



$m(4,3)=20: 20$-III. в $P(4,3)$ не единственна;



$m(5,3)=56: 56$-III. в $P(5,3)$ единственна.



III. используются в теории кодирования (см., например, [2]).



[2] Jum.: [1] B o s e R. C., «Sankhyă», 1947, v. 8, p. 107-66; [2] H i l l R., "Discrete Math.", 1978, v. 22, №2, p. 111-37; p. $133-236$.

B. B. Аøанасьев.



III AP -- множество $V^{n}$ точек $x$ евкјидова пространства $E^{n}$, удаленных от нек-рой точки $x_{0}$ (цент р III.) на расстояние, меньшее (о т к р т ы й ш а р $\dot{V}^{n}$ ) или не превышаюшее ( а м к н у т ы й ш а р $\bar{V} n$ ) величину $R$ (р а д и у с Ш.), т. е.



\[

\stackrel{\circ}{V}\left(\overline{V^{n}}\right)=\left\{x \in E^{n}, \rho\left(x, x_{0}\right)<R(\leqslant R)\right\} .

\]



III. $V^{1}$ - это о т р е зок, $V^{2}$ - эток р у г, $V^{n}$ при $n>3$ иногда наз. г и п е р ш а о м. Границей (поверхностью) ШІ. является сфера.



Объем Ill.



\[

V^{n}=\frac{\pi^{n / 2}}{\Gamma(n / 2+1)} R^{n}=\frac{S^{n-1}}{n} R,

\]



где $S^{n-1}$ - поверхность граничной сферы. В частності,



\[

V^{1}=2 R, V^{2}=\pi R^{2}, V^{3}=\frac{4}{3} \pi R^{3}, V^{4}=\frac{\pi^{2} R^{4}}{2} .

\]



С возрастанием размерности объем ІІІ. «концентрируется» у его поверхности:



\[

\frac{V^{n}\left(\boldsymbol{R}_{2}\right)-V^{n}\left(\boldsymbol{R}_{1}\right)}{V^{n}\left(\boldsymbol{R}_{2}\right)} \longrightarrow 1, R_{1}<R_{2}, n \longrightarrow \infty .

\]



Ш.- простейшая геометрич. фигура. Топология его тривиальна. Среди всех тел равного объема III. имеет минимальную поверхность, а среди всех тел одинаковой поверхности он имеет наименьший объем.



Совершенно аналогично определяется III. в метрич. пространстве, однако при этом, напр., ІІІ. может быть не строго выпуклым, его поверхность может иметь негладкие точки и т. п.- все то, что бывает у произвольных выпуклых тел.



Бе с к о н е ч в о ме рвы й ші а $\mathrm{p}-$ прямой предел последовательности вложенных друг в друга III. последовательных размерностей - в отличие от конечномерного ШII. не имеет компактного замыкания. Наоборот, компантность ШІ. в топологическом векторном пространстве свидетельствует о его конечномерности.



Лит. см. при статье Сøера. Н. С. Шарадзе.



IIIАРЛЬЕ МНОГОЧЛЕНЫ - многочленЫ, ортогональные на системе неотрицательных целотисленных точек с интегральным весом $d \sigma(x)$, где $\sigma(x)$ - ступенчатая функция, скачки к-рой определяются по формуле



\[

j(x)=e^{-a} \frac{a^{x}}{x !}, x=0,1,2, \ldots ; a>0 .

\]



Ортонормированные Ш. м. имеют представления



\[

\begin{aligned}

P_{n}(x ; a) & =\sqrt{\frac{a^{n}}{n !}} \sum_{k=0}^{n}(-1)^{n-k}\left(\begin{array}{l}

n \\

k

\end{array}\right) k ! a^{-k}\left(\begin{array}{l}

x \\

k

\end{array}\right)= \\

& =a^{\frac{n}{2}}(n !)^{-\frac{1}{2}}[j(x)]^{-1} \Delta_{j}^{n}(x-n) .

\end{aligned}

\]



С Лагерра многочленами ШІ. м. связаны равенством



\[

P_{n}(x ; a)=\sqrt{\frac{n !}{a^{n}}} L_{n}^{(x-n)}(a)=\sqrt{\frac{n !}{a^{n}}} L_{n}(a ; x-n) .

\]



Введены К. Шарлье [1]. Поскольку функция $j(x)$ определяет распределение Пуассона, то многочлены $\left\{P_{n}(x ; a)\right\}$ наз. также многочленами П у ас с она а-II а р лье.



Jum.: [1] C harlie r C., Application de la théorie des probabilités à l'astronomie, P., 1931; [2] Б е й т м е н $\mathbf{\Gamma}$., Ә pд е й и А., Высшие трансцендентные функции. Функции Бессе-

ля, функции параболического цилиндра, ортогональные многоля, функции параболического цилиндра, ортогональные много-

члены, пер. с англ., М., 1974; [3] С е г ё Г., Ортогональные многочлены, пер. с анг,



IIАРЛЬЕ РАСПРЕДЕЛЕНИЕ - малоупотребительное название распределения, плотность к-рого дается Грама - UІарлье рлдом.



IIIАРОВАЯ ФУНКЦИЯ, т е л е с н а я г а р м о н и ч еская ф у нкци я,- сферическая функция п-й степени с множителем $r^{n}$.



IIАУДЕРА МЕТОД - метод решения краевых задач для линейных равномерно эллиптических уравнений 2-го ıорядка, в основе к-рого лежкат априорные оценки и метод продолжения по параметру.



II. м. решения Дирихле задачи для линейного равномерно эллинтического уравнения



\[

L u \equiv a^{i j}(x) u_{x_{i} x_{j}}+b^{j}(x) u_{x_{j}}+b(x) u=f(x),

\]





\end{document}